\documentclass[11pt,a4paper]{article}

\usepackage[T1]{fontenc}
\usepackage[utf8]{inputenc}
\usepackage{lmodern}
\usepackage{microtype}
\usepackage[a4paper,margin=1in]{geometry}
\usepackage{amsmath,amssymb,mathtools,bm}
\usepackage{booktabs}
\usepackage[hidelinks]{hyperref}

\newcommand{\R}{\mathbb{R}}
\newcommand{\dd}{\,\mathrm{d}}
\newcommand{\bn}{\bm{n}}
\newcommand{\bx}{\bm{x}}
\newcommand{\by}{\bm{y}}
\newcommand{\bX}{\bm{X}}
\newcommand{\bu}{\bm{u}}
\newcommand{\bU}{\bm{U}}
\newcommand{\bomega}{\bm{\omega}}
\newcommand{\bOmega}{\bm{\Omega}}
\newcommand{\bL}{\bm{L}}
\newcommand{\bLam}{\bm{\Lambda}}
\newcommand{\bF}{\bm{F}}
\newcommand{\bG}{\bm{G}}
\newcommand{\bp}{\bm{p}}
\newcommand{\bh}{\bm{h}}
\newcommand{\bI}{\bm{I}}
\newcommand{\bM}{\bm{M}}
\newcommand{\bC}{\bm{C}}
\newcommand{\bz}{\bm{z}}
\newcommand{\bd}{\bm{d}}
\newcommand{\bc}{\bm{c}}
\newcommand{\bA}{\bm{A}}
\newcommand{\bK}{\bm{K}}
\newcommand{\bq}{\bm{q}}
\newcommand{\quat}[1]{\mathbf{#1}}

\title{Equations Used by the Multi-Ellipsoid Boundary-Integral Solver}
\author{}
\date{\today}

\begin{document}
\maketitle

\begin{abstract}
This note records the mathematical model and discrete equations implemented in
the current multi-ellipsoid solver.  The fluid is inviscid, incompressible and
irrotational, so the hydrodynamic coupling is obtained from a Neumann
boundary-integral solve for the velocity potential.  The preferred multibody
time step is an impulse/Hamiltonian endpoint formulation: fluid impulse and
total kinetic energy are treated as state functions of the rigid-body
configuration and velocities, and the endpoint velocities are chosen to satisfy
the global linear impulse, angular impulse and energy constraints.
\end{abstract}

\tableofcontents

\section{Rigid bodies and notation}

The system consists of $N$ rigid ellipsoids immersed in an unbounded ideal
fluid of density $\rho_f$.  Body $b$ has centre $\bX_b$, orientation quaternion
$\quat{q}_b$, semi-axes $(a_b,b_b,c_b)$, solid density $\rho_{s,b}$ and volume
\begin{equation}
  V_b = \frac{4\pi}{3}a_b b_b c_b .
\end{equation}
The solid mass and body-frame inertia tensor are
\begin{equation}
  m_b = \rho_{s,b} V_b,\qquad
  \bI_{s,b} =
  \frac{m_b}{5}
  \begin{bmatrix}
    b_b^2+c_b^2 & 0 & 0\\
    0 & a_b^2+c_b^2 & 0\\
    0 & 0 & a_b^2+b_b^2
  \end{bmatrix}.
  \label{eq:solid-mass-inertia}
\end{equation}
The lab-frame linear velocity is $\bU_b=\dot{\bX}_b$.  The lab angular velocity
is represented as the pure quaternion $(0,\bOmega_b)$; body-frame vectors are
obtained by
\begin{equation}
  (0,\bomega_b) = \quat{q}_b^{-1}(0,\bOmega_b)\quat{q}_b .
  \label{eq:lab-body-rotation}
\end{equation}
The solid angular momentum in the lab frame is therefore
\begin{equation}
  \bh_{s,b}
  = \quat{q}_b\,(0,\bI_{s,b}\bomega_b)\,\quat{q}_b^{-1}.
  \label{eq:solid-angular-momentum}
\end{equation}
The solid kinetic energy is
\begin{equation}
  T_{s,b}
  = \frac{1}{2}m_b \|\bU_b\|^2
  + \frac{1}{2}\bomega_b^\mathsf{T}\bI_{s,b}\bomega_b .
  \label{eq:solid-ke}
\end{equation}

\section{Single-ellipsoid added-mass tensors}

For an isolated ellipsoid the solver also computes the classical added-mass
reference tensors from the shape functions
\begin{equation}
  \alpha_i
  = abc\int_0^\lambda
  \frac{\dd s}{(a_i^2+s)
  \sqrt{(a^2+s)(b^2+s)(c^2+s)}} ,
  \qquad (a_1,a_2,a_3)=(a,b,c).
  \label{eq:shape-functions}
\end{equation}
In the code this integral is evaluated with a large finite upper limit
$\lambda=10^4$.  Writing $(\alpha_1,\alpha_2,\alpha_3)=(\alpha,\beta,\gamma)$,
the translational added-mass tensor is
\begin{equation}
  \bM_f
  = \rho_f V
  \operatorname{diag}\left(
    \frac{\alpha}{2-\alpha},
    \frac{\beta}{2-\beta},
    \frac{\gamma}{2-\gamma}
  \right).
  \label{eq:added-mass}
\end{equation}
The rotational fluid inertia reference used for diagnostics is
\begin{equation}
  \bI_f
  = \frac{\rho_f V}{5}
  \operatorname{diag}(e_1,e_2,e_3),
  \label{eq:fluid-inertia}
\end{equation}
where
\begin{align}
 e_1 &=
 \frac{(b^2-c^2)^2(\gamma-\beta)}
 {2(b^2-c^2)+(\beta-\gamma)(b^2+c^2)},\\
 e_2 &=
 \frac{(a^2-c^2)^2(\gamma-\alpha)}
 {2(a^2-c^2)+(\alpha-\gamma)(a^2+c^2)},\\
 e_3 &=
 \frac{(a^2-b^2)^2(\beta-\alpha)}
 {2(a^2-b^2)+(\alpha-\beta)(a^2+b^2)}.
\end{align}
For a single isolated body the reference Hamiltonian is then
\begin{equation}
  H_{\mathrm{iso}}
  = \frac{1}{2}\bU^\mathsf{T}(\bM_s+\bM_f)\bU
  + \frac{1}{2}\bomega^\mathsf{T}(\bI_s+\bI_f)\bomega .
  \label{eq:single-body-reference-hamiltonian}
\end{equation}

\section{Potential-flow boundary integral}

The fluid velocity is $\bu=\nabla\phi$ and the potential satisfies
\begin{equation}
  \nabla^2\phi = 0
  \quad\text{in the exterior fluid domain}, \qquad
  \nabla\phi(\bx)\to 0\quad\text{as }\|\bx\|\to\infty .
  \label{eq:laplace}
\end{equation}
On the surface $S_b$ of body $b$, the no-penetration condition gives the
Neumann data
\begin{equation}
  \frac{\partial\phi}{\partial n}(\bx)
  = \bn(\bx)\cdot
  \left[\bU_b+\bOmega_b\times(\bx-\bX_b)\right],
  \qquad \bx\in S_b .
  \label{eq:neumann}
\end{equation}
The free-space Green function and its gradient are
\begin{equation}
  G(\bx,\by)=\frac{1}{4\pi\|\bx-\by\|},\qquad
  \nabla_{\bx} G(\bx,\by)
  =-\frac{\bx-\by}{4\pi\|\bx-\by\|^3}.
  \label{eq:green}
\end{equation}
The direct boundary-integral equation is assembled in the form
\begin{equation}
  c(\bx)\phi(\bx)
  + \int_S \phi(\by)
    \frac{\partial G(\bx,\by)}{\partial n_{\by}}\,\dd S_{\by}
  =
  \int_S G(\bx,\by)
    \frac{\partial\phi}{\partial n}(\by)\,\dd S_{\by} ,
  \qquad \bx\in S .
  \label{eq:bem-continuous}
\end{equation}
After discretisation on quadratic triangular elements this becomes the dense
linear system
\begin{equation}
  \bA(\bX,\quat{q})\,\bm{\phi}
  =
  \bm{b}(\bX,\quat{q},\bU,\bOmega),
  \label{eq:bem-discrete}
\end{equation}
with $\bm{\phi}$ the nodal boundary potential.  Matrix entries are generated by
double-layer quadrature and the right-hand side by the single-layer integral of
the imposed Neumann data.  In the current multibody solve all bodies share one
combined surface mesh, so cross-body hydrodynamic interactions are contained in
the off-diagonal blocks of $\bA$.

\section{Surface state functions}

Once \eqref{eq:bem-discrete} has been solved, the fluid kinetic energy is
computed by Green's identity as the surface integral
\begin{equation}
  T_f
  =
  -\frac{\rho_f}{2}\int_S
  \phi\,\frac{\partial\phi}{\partial n}\,\dd S .
  \label{eq:fluid-ke}
\end{equation}
The minus sign follows the outward-normal convention of the exterior fluid
problem used by the implementation.

The per-body Lamb impulse and angular Lamb impulse are
\begin{align}
  \bL_b
  &= \rho_f\int_{S_b}\phi\,\bn\,\dd S,\\
  \bLam_b
  &= \rho_f\int_{S_b}
  \phi\,(\bx-\bX_b)\times\bn\,\dd S .
  \label{eq:lamb-impulse}
\end{align}
These are state functions of the instantaneous configuration and generalized
velocity.  They are central to the impulse schemes because they do not require a
history-dependent approximation to $\partial\phi/\partial t$.

For the pressure-force option, the unsteady potential term gives
\begin{align}
  \bF_b^{(\dot\phi)}
  &= \rho_f\int_{S_b}\dot\phi\,\bn\,\dd S,\\
  \bm{N}_b^{(\dot\phi)}
  &= \rho_f\int_{S_b}
  \dot\phi\,(\bx-\bX_b)\times\bn\,\dd S.
  \label{eq:phidot-force}
\end{align}
When the surface is rotating, the full time derivative of the impulse also
contains transport terms.  The implemented pressure-mode correction is
\begin{align}
  \bF_b &= \bF_b^{(\dot\phi)} + \bOmega_b\times\bL_b,\\
  \bm{N}_b &= \bm{N}_b^{(\dot\phi)} + \bOmega_b\times\bLam_b .
  \label{eq:transport}
\end{align}
The missing angular version of \eqref{eq:transport} was the source of the old
single-body impulse energy error.

\section{Conserved impulse and Hamiltonian}

The solver monitors and, in the Hamiltonian endpoint scheme, enforces the total
linear and angular conserved impulses
\begin{align}
  \bm{P}
  &=
  \sum_{b=1}^N
  \left(m_b\bU_b-\bL_b\right),
  \label{eq:linear-impulse}\\
  \bm{H}
  &=
  \sum_{b=1}^N
  \left[
    \bX_b\times\left(m_b\bU_b-\bL_b\right)
    + \bh_{s,b}
    - \bLam_b
  \right].
  \label{eq:angular-impulse}
\end{align}
The total kinetic energy used by the energy-conserving schemes is
\begin{equation}
  H(\bX,\quat{q},\bU,\bOmega)
  =
  \sum_{b=1}^N T_{s,b}
  + T_f .
  \label{eq:total-hamiltonian}
\end{equation}
Equations \eqref{eq:linear-impulse}--\eqref{eq:total-hamiltonian} are evaluated
from the same BEM solve, so all added-mass and body--body interaction terms are
included implicitly.

\section{Impulse-difference time step}

Let $h$ be the time step and let $(\bL_b^n,\bLam_b^n)$ denote the BEM impulse
at the start of the step.  The implicit impulse-difference update solves for
the endpoint velocity and a step-averaged torque.  The translational equation is
\begin{equation}
  m_b\left(\bU_b^{n+1}-\bU_b^n\right)
  =
  \bL_b^{n+1}-\bL_b^n
  + h\,\bG_{X,b},
  \label{eq:impulse-linear-step}
\end{equation}
where $\bG_X$ is normally zero and is only nonzero when the experimental finite
difference gradient of the fluid kinetic energy is enabled.  Positions use the
midpoint velocity,
\begin{equation}
  \bX_b^{n+1}
  =
  \bX_b^n+\frac{h}{2}
  \left(\bU_b^n+\bU_b^{n+1}\right).
  \label{eq:impulse-position}
\end{equation}
The consistent torque uses the angular impulse difference about the translating
body centre and the reference-point transport correction
\begin{equation}
  \bm{N}_b
  =
  \frac{\bLam_b^{n+1}-\bLam_b^n}{h}
  - \bar{\bU}_b\times\bar{\bp}_b
  + \bm{G}_{q,b},
  \qquad
  \bar{\bp}_b
  =
  \frac{1}{2}
  \left[
  m_b\bU_b^n-\bL_b^n
  +m_b\bU_b^{n+1}-\bL_b^{n+1}
  \right].
  \label{eq:impulse-torque-step}
\end{equation}
The angular velocity is advanced with a Lie-group implicit midpoint rule in the
body frame:
\begin{align}
  \bomega_{b,m}
  &=
  \bomega_b^n
  + \frac{h}{2}\bI_{s,b}^{-1}
  \left[
    \bm{N}_{b,m}
    -\bomega_{b,m}\times
    \left(\bI_{s,b}\bomega_{b,m}\right)
  \right],
  \label{eq:angular-midpoint}\\
  \quat{q}_b^{n+1}
  &=
  \exp\left(\frac{h}{2}(0,\bOmega_{b,m})\right)\quat{q}_b^n,
  \label{eq:quat-step}\\
  \bomega_b^{n+1}
  &=2\bomega_{b,m}-\bomega_b^n .
  \label{eq:omega-end}
\end{align}
Here $\bm{N}_{b,m}$ is the applied torque rotated into the midpoint body frame.
The nonlinear fixed point stacks the unknown endpoint velocities and torques,
and Anderson acceleration is applied to that stacked residual.

\section{Hamiltonian midpoint endpoint solve}

The most conservative multibody option solves directly for generalized
endpoint data.  Define the generalized velocity vector
\begin{equation}
  \bz =
  \left[
  \bU_1,\ldots,\bU_N,
  \bOmega_1,\ldots,\bOmega_N
  \right]^\mathsf{T}\in\R^{6N}.
\end{equation}
For a midpoint-velocity version of the method, $\bz$ defines the endpoint
configuration by
\begin{align}
  \bX_b^{n+1}(\bz)
  &= \bX_b^n + h\,\bU_{b,m},\\
  \quat{q}_b^{n+1}(\bz)
  &= \exp\left(\frac{h}{2}(0,\bOmega_{b,m})\right)\quat{q}_b^n .
  \label{eq:endpoint-from-mid}
\end{align}
The midpoint configuration used for the BEM solve is reconstructed from the
start and endpoint states:
\begin{align}
  \bX_{b,m}
  &=\frac{1}{2}\left(\bX_b^n+\bX_b^{n+1}\right),\\
  \bOmega_{b,m}
  &= \frac{1}{h}\log\left(\quat{q}_b^{n+1}(\quat{q}_b^n)^{-1}\right),\\
  \quat{q}_{b,m}
  &= \exp\left(\frac{h}{4}(0,\bOmega_{b,m})\right)\quat{q}_b^n .
  \label{eq:midpoint-reconstruction}
\end{align}
The endpoint/coupled variant instead evaluates the conservation residual at
the endpoint state generated by $\bz$.

Let
\begin{equation}
  \bC(\bz)=
  \begin{bmatrix}
    \bm{P}(\bz)\\
    \bm{H}(\bz)
  \end{bmatrix},
  \qquad
  \bC^\star=\bC(\bz^n),
  \qquad
  H^\star=H(\bz^0),
  \label{eq:target-invariants}
\end{equation}
where $H^\star$ is the initial total kinetic energy.  The seven-component
nonlinear residual used by the coupled endpoint solve is
\begin{equation}
  \bm{R}(\bz)=
  \begin{bmatrix}
  \left(\bm{P}(\bz)-\bm{P}^\star\right)/s_P\\
  \left(\bm{H}(\bz)-\bm{H}^\star\right)/s_H\\
  \left(H(\bz)-H^\star\right)/s_E
  \end{bmatrix},
  \label{eq:coupled-residual}
\end{equation}
with scales
\begin{equation}
  s_P=\max(\|\bm{P}^\star\|,1),\qquad
  s_H=\max(\|\bm{H}^\star\|,1),\qquad
  s_E=\max(|H^\star|,1).
\end{equation}
The code solves \eqref{eq:coupled-residual} by finite-difference Newton steps,
with optional Broyden rank-one Jacobian updates and a trust-style cap on the
generalized velocity increment.

\section{Energy and impulse projection}

The projection used by the Hamiltonian experiments is algebraic.  At fixed
configuration the conserved impulse is affine in generalized velocity,
\begin{equation}
  \bC(\bz)=\bA\bz+\bc_0,
  \label{eq:affine-impulse}
\end{equation}
where $\bA$ is obtained by evaluating the BEM state function on basis velocity
vectors.  A first correction projects the current velocity $\bz_c$ onto the
linear impulse constraint:
\begin{equation}
  \bz_P
  =
  \bz_c
  + \bA^\mathsf{T}(\bA\bA^\mathsf{T})^{-1}
  \left(\bC^\star-\bC(\bz_c)\right).
  \label{eq:linear-projection}
\end{equation}
A minimum-energy particular solution satisfying the same affine constraint is
computed from the sampled energy quadratic
\begin{equation}
  H(\bz)=H(\bm{0})+\frac{1}{2}\bz^\mathsf{T}\bK\bz ,
  \label{eq:energy-quadratic}
\end{equation}
as
\begin{equation}
  \bz_0
  =
  \bK^{-1}\bA^\mathsf{T}
  \left(\bA\bK^{-1}\bA^\mathsf{T}\right)^{-1}
  \left(\bC^\star-\bc_0\right).
  \label{eq:min-energy-particular}
\end{equation}
The remaining correction is restricted to the null direction
\begin{equation}
  \bd = \bz_P-\bz_0,\qquad \bz(\alpha)=\bz_0+\alpha\bd .
\end{equation}
The scalar $\alpha$ is chosen from
\begin{equation}
  H(\bz_0+\alpha\bd)=H^\star .
  \label{eq:scalar-energy-projection}
\end{equation}
In implementation the coefficients of this quadratic are evaluated by three
BEM energy calls:
\begin{align}
  q_a &= \frac{1}{2}\left[H(\bz_P)+H(\bz_0-\bd)\right]-H(\bz_0),\\
  q_b &= \frac{1}{2}\left[H(\bz_P)-H(\bz_0-\bd)\right],\\
  q_c &= H(\bz_0)-H^\star,
\end{align}
and the root closest to $\alpha=1$ is selected.  If the energy floor
$H(\bz_0)$ already exceeds the target, the exact energy projection is
infeasible and the code falls back to $\bz_0$ while reporting the floor error.

\section{Diagnostics written to output}

The CSV output writes the total energy components
\begin{equation}
  H,\qquad T_f,\qquad \sum_b T_{s,b},
\end{equation}
and, for each body, the diagnostic vectors
\begin{align}
  \bp_{\mathrm{con},b} &= m_b\bU_b-\bL_b,\\
  \bh_{\mathrm{con},b} &=
  \bX_b\times\bp_{\mathrm{con},b}
  +\bh_{s,b}-\bLam_b .
\end{align}
For a closed multibody system the sums of these vectors are the conserved
quantities in \eqref{eq:linear-impulse} and \eqref{eq:angular-impulse}.  The
study dashboards report kinetic-energy drift, separation, trajectories,
orientation markers, and relative drift in the summed conserved linear and
angular impulses.

\section{Remarks on physical interpretation}

The BEM solve makes the fluid energy and impulse depend on the entire multibody
configuration, not on a pairwise approximation.  Low-density bodies therefore
have small solid inertia relative to the configuration-dependent fluid inertia,
which is precisely the regime in which explicit pressure-force updates drift
or destabilise.  The impulse/Hamiltonian schemes instead enforce the global
invariants \eqref{eq:linear-impulse}--\eqref{eq:total-hamiltonian} at the
endpoint level, making the energy behaviour much closer to the conservative
ideal-fluid dynamics.

\end{document}
